% ECONOMICS_PROBLEM: {"id": "D91-431", "title": "Causal mental models", "jel_code": "D91", "source_id": "OP-0304"}
% !TeX program = xelatex
\documentclass[11pt,a4paper]{article}
\usepackage[margin=23mm]{geometry}
\usepackage{fontspec}
\setmainfont{Latin Modern Roman}
\usepackage{amsmath,amssymb,mathtools}
\usepackage{xcolor,enumitem,booktabs,longtable,array,xurl,hyperref}
\definecolor{muted}{HTML}{53636C}
\hypersetup{colorlinks=true,urlcolor=blue,linkcolor=blue}
\setlength{\parindent}{0pt}
\setlength{\parskip}{5pt}
\newcommand{\R}{\mathbb R}
\newcommand{\E}{\mathbb E}
\newcommand{\N}{\mathbb N}
\newcommand{\ind}{\mathbf 1}
\newcommand{\Var}{\operatorname{Var}}
\newcommand{\Cov}{\operatorname{Cov}}
\newcommand{\supp}{\operatorname{supp}}
\DeclareMathOperator*{\argmin}{arg\,min}
\DeclareMathOperator*{\argmax}{arg\,max}
\newcommand{\entrymeta}[3]{{\sffamily\small\color{muted}\textbf{\texttt{#1}}\quad|\quad #2\par #3\par}}
\newcommand{\Problemref}[1]{\texttt{#1}}
\newcommand{\JELref}[2]{\texttt{#2} (JEL \texttt{#1})}
\begin{document}
\section*{Causal mental models}
\textbf{Problem ID:} \texttt{D91-431}\quad \textbf{JEL:} \texttt{D91}\par
% BEGIN ECONOMICS STATEMENT
\label{problem:OP-0304}
\label{audited:ECON-023}
{\sffamily\small\color{muted}Original subject group: Behavioral and experimental economics\par}
\label{definitions:problem:30007586}
\textbf{Audit classification:} Published research agenda. The genuine 2025 review §§2–3 and §5 discusses rationalization, elicitation reactivity and incentives. A DAG reconstruction from text is an editorial statistical specification.
\subsection*{Definitions and precise research target}
\textbf{Model and research target.} Let individual \(i\)'s latent reasoning process be \(M_i\), represented by a directed acyclic graph over specified economic variables and a set of perceived conditional distributions. A task presents information \(X\), and the individual chooses \(A_i\). An elicitation protocol \(E\) produces text \(T_i\), which a fixed coding rule \(h\) converts into reported model \(\widehat M_i=h(T_i)\). Crucially, both \(A_i(E)\) and \(T_i(E)\) may depend on elicitation timing, incentives, and prompts. Consider randomized counterfactual-information interventions \(Z\) chosen to distinguish models in a prespecified finite class \(\mathcal M\). Define model predictive loss and elicitation reactivity as
\[R=\mathbb E[\ell(A_i(Z),p_{\widehat M_i}(Z))],\qquad \rho=\mathbb E[A_i(E_1)-A_i(E_0)].\]
Here \(E_0\) omits explanatory prompts and \(E_1\) elicits explanations before or during choice, with all other task features fixed. Identify conditions under which observed text and intervention responses distinguish genuine pre-choice reasoning from retrospective justification, while estimating or bounding \(\rho\). High agreement between coders identifies reproducible classification, not necessarily latent-model accuracy. Multiple latent processes may yield the same text and choices, so point identification requires exclusions, informative interventions, or independent process measurements. The empirical objective is to establish validation criteria and uncertainty bounds for recovered models without assuming that verbal reports transparently reveal the causal mechanisms governing behavior.

\emph{Definitions and maintained assumptions (editorial unless a source definition is named).} Fix variable set \(V\), finite DAG class \(\mathcal M\), and a perceived probability model \(P_M\) satisfying the DAG’s factorization; a graph alone does not define predicted action. Specify how observations and interventions enter the subjective model, the choice utility and the resulting response law. A fixed coding map \(h\) uses a prespecified rubric or frozen classifier and outputs a graph, ambiguity set or unclassifiable category. Validation loss must be operational: for instance \(V=\mathbb E[\ell(A_i(Z),p_{\widehat M_i}(Z))]\) with a declared proper loss, rather than the undefined event ``predicts.'' Elicitation effects compare randomized protocols \(E_1,E_0\) holding all other inputs fixed, giving \(\rho=\mathbb E[A(E_1)-A(E_0)]\) for scalar actions. Classifier agreement measures coding reproducibility. Latent-model identification requires distinguishing response laws across informative interventions and a measurement model \(P(T\mid M,E)\); retrospective narratives may be generated by a different latent process from pre-choice reasoning.
\subsection*{Variants and assumption changes}
\begin{enumerate}
\item \textbf{Timing validation (editorial).} Randomize pre-choice versus post-choice explanation requests and no-explanation controls. Estimate action reactivity and held-out predictive losses, preserving a category for incomparable or unclassifiable narratives.
\item \textbf{Intervention identification (editorial).} For a prespecified finite perceived-DAG class, choose interventions on variables that distinguish model predictions. Characterize remaining equivalence classes when text-report error or strategic self-presentation is admitted.
\end{enumerate}
\subsection*{References and source audit}
\begin{enumerate}
\item §3.1.1 p.14; §5 pp.40–41 in April2025 manuscript. Supports the stated research area; exact displayed specification is editorial. \url{https://www.econtribute.de/RePEc/ajk/ajkdps/ECONtribute_362_2025.pdf}
\end{enumerate}
\begin{enumerate}\raggedright
\item\hypertarget{definitions:source:30007586:1}{} Ingar Haaland; Christopher Roth; Stefanie Stantcheva; Johannes Wohlfart. 2025. \emph{Understanding Economic Behavior Using Open-Ended Survey Data}. §3.1.1 p.14; §5 pp.40–41 in April2025 manuscript.
\url{https://www.econtribute.de/RePEc/ajk/ajkdps/ECONtribute_362_2025.pdf}.
DOI: \url{https://doi.org/10.1257/jel.20251780}.
\end{enumerate}

\subsection*{Integrated refinement: ECON-023}
{\sffamily\small\color{muted}Belief causation versus narrative rationalization\par}
This refinement adopts the additional source conventions
(page~\pageref{audited:conventions}), subject to its local restrictions.
\textbf{Belief mediation versus retrospective rationalization.}
With randomized information $Z_i$, pre-choice beliefs $B_i(z)$ and choices
$Y_i(z,b)$, distinguish narratives generated by beliefs from narratives generated
after choices. A specific extension is
\[
\theta_B=\mathbb E[Y_i(0,B_i(1))-Y_i(0,B_i(0))].
\]
Specify timing, exclusion, the belief/text measurement law, and cross-world
restrictions identifying or bounding this contrast. Information randomization
identifies total effects under its assignment assumptions; it does not by itself
identify the mediated contrast. Randomize elicitation timing and preserve
pre-choice, post-choice and no-explanation controls. Distinguish interventions on
beliefs from changes in preferences or opportunities.
\subsection*{Supplied source audit for ECON-023}
Local [S1], [S2], etc. in this audit and the references immediately below
belong to ECON-023, independently of the original entry's reference numbering.
[S1]'s working-paper introduction distinguishes decision-relevant reasoning from rationalization. Its discussion supports further work linking beliefs and choices. This natural indirect effect is editorial and is not claimed to be identified by the survey.
\subsection*{References for integrated refinement ECON-023}
{\footnotesize\raggedright
\par\noindent\textbf{[S1]} Ingar Haaland, Christopher Roth, Stefanie Stantcheva, and Johannes Wohlfart (2025). \emph{Understanding Economic Behavior Using Open-ended Survey Data}. Journal of Economic Literature 63(4), 1244–1280. DOI: 10.1257/jel.20251780. Page locators below refer to the April 2025 ECONtribute Discussion Paper 362 version.\par \textit{Verified locator / source role:} Introduction, printed p. 4 (PDF page 6): genuine mechanisms versus ex-post rationalization; Section 5, Neuroeconomics, p. 41.\par \url{https://www.econtribute.de/RePEc/ajk/ajkdps/ECONtribute_362_2025.pdf}\par\smallskip
}

\subsection*{Common conventions: Source mathematical conventions}

These conventions apply to each entry unless explicitly replaced.

\textbf{Models and identification.} A primitive model \(m\) specifies all probability laws, feasible choices, information, equilibrium and measurement rules used by its target. The admissible class \(\mathcal M\) is fixed independently of outcomes. If \(P_O(m)\) is its observable probability law and \(\tau(m)\) the target, its identified set at observable law \(P\) is
\[
 \mathcal I_\tau(P)=\{\tau(m):m\in\mathcal M,\ P_O(m)=P\}.
\]
Point identification means that this set is a singleton. An empty set signals incompatibility of data and assumptions. Coordinate bounds need not describe the joint set. Bounds are sharp only if every retained target is attainable or arbitrarily approachable by compatible models. Finite bounds require explicit restrictions on unobserved quantities; unrestricted confounding, hidden wealth, extrapolation or arbitrary equilibrium selection can yield uninformative sets.

\textbf{Causal interventions.} A potential outcome \(Y_i(p)\) is the outcome under a complete policy regime \(p\), including timing, financing, information and market spillovers. The target population and units are fixed before treatment. With interference, use the complete assignment vector or a justified exposure mapping; individual no-interference notation alone cannot identify an economy-wide intervention. A difference of mean potential outcomes does not require the cross-world joint distribution. Conditioning on post-treatment migration, survival, enrollment or employment changes the estimand and needs separate justification.

\textbf{Statistical statements.} An estimator, confidence set, forecast or test specifies a sampling law, model class and loss/coverage criterion. Uniform \((1-\alpha)\)-coverage means \(\inf_{m\in\mathcal M}\Pr_m\{\tau(m)\in C_n\}\geq1-\alpha\), or an explicitly asymptotic version. The whole parameter space trivially has coverage, so informativeness requires a stated width, risk or power objective. A forecasting improvement is relative to a named benchmark, information set and evaluation population.

\textbf{Equilibrium and optimization.} An equilibrium comprises feasible plans, optimality, beliefs where needed, budget constraints and all market/resource clearing. The primitives or a stated benchmark define each correspondence \(\mathcal E(p;m)\). Nonemptiness is assumed or proved, never inferred just from compact policy sets. A selected-equilibrium target uses a declared selection rule; a robust target quantifies over all permitted equilibria. Use suprema or infima unless attainment is established. An \(\varepsilon\)-optimal policy is feasible and within \(\varepsilon\) of the finite optimum value. Report an empty feasible set separately.

\textbf{Welfare and accounting.} Normative utility, interpersonal normalization, social weights, numeraire, discounting and the use of public receipts are primitives. Behavioral utility need not coincide with the welfare criterion. Household welfare includes ownership income and rents once; adding profits again to compensating variation generally double-counts them. A monetary equivalent is defined by an explicit transfer and price convention, with monotonicity and range conditions. Average effects, mechanism contributions, transfers and welfare losses are different targets. Infinite sums require convergence and dynamic feasible plans require terminal or no-Ponzi conditions.

\textbf{Source versions.} A working-paper date, revision date and journal publication year are distinguished. A DOI for a journal version is not automatically the DOI of an earlier manuscript. Locators refer to the stated version and printed pages where specified. A metadata-only check does not verify a claim about a full-text conclusion. Every entry's source subsection narrows the provenance claim accordingly.

\subsection*{Common conventions: Additional-source status and mathematical conventions}

\label{audited:conventions}

Of the 100 supplied ECON entries, 65 distinct problems appear here; 32 close
matches refine earlier entries, and three duplicate questions map to existing
entries. Appendix~\ref{app:audited-crosswalk} lists every destination.
The attachment's P/R/S/U distinctions and source audits are retained. Its
precise mathematical formalizations are editorial unless the individual audit
says otherwise. Candidate subsidy targets ECON-004--006 retain their U flags;
the resolved approval-core question ECON-007 updates OP-0202 above.
The following supplied conventions govern both this part and the attributed
ECON refinements elsewhere in the catalogue, subject to local restrictions.

\textbf{Parameterized questions.} A problem family lists inputs that must be fixed before an instance is evaluated: population, horizon, policies, information, data law, model class and welfare weights. These are required choices, not quantities the citations are assumed to specify. Undefined applications should not be treated as identified estimands.

\textbf{Indices and integrability.} Sets of agents, firms and regions are finite unless a population measure is explicitly used. \(i,j\) index units, \(r\) regions, \(t\) dates and \(h\) horizons. \(\mathbb E\) and \(\Pr\) refer to the declared population and law; \(\mathbf1\{A\}\) is an indicator. Every displayed expectation and discounted sum needs existence in the stated domain. Infinite horizons use a discount in \((0,1)\) and a convergence condition; finite horizons may use discount one.

\textbf{Optimization and existence.} A displayed \(\max\), \(\min\), or \(\arg\max\) is conditional on attainment. For finite feasible menus this is immediate; general spaces need continuity and compactness/coercivity or another existence argument. Without attainment, the target is the supremum/infimum and arbitrarily near-optimal feasible choices. \(\inf\varnothing=+\infty\). Empty feasibility and an unidentified target are different issues.

\textbf{Potential outcomes.} \(Y_i(z)\) is the outcome under a specified intervention, not the observed conditional mean among people choosing \(z\). In binary comparisons \(\Delta Y_i=Y_i(1)-Y_i(0)\). Specify assignment, treatment versions, compliance, information and observation horizon. Interference is allowed under a full policy vector; compressed exposure notation needs a justified mapping. No interference, consistency and overlap are substantive assumptions, not automatic consequences of notation.

\textbf{Populations and observation.} Fix baseline eligibility and population weights. Conditioning on post-policy employment, survival, relocation or program uptake can change the population being compared. Death is not ordinary loss to follow-up; outcomes truncated by death need a composite convention, joint survival target, or explicitly defined survivor stratum. Fertility-changing interventions need a population functional rather than an unexplained fixed descendant cohort. Measurement and missingness models must be supplied.

\textbf{Identification.} A structural model \(m\in\mathcal M\) implies observable law \(P_m\) and target \(\theta(m)\). Identification means
\[
 P_{m_1}=P_{m_2}\ \Longrightarrow\ \theta(m_1)=\theta(m_2)
 \quad\text{for all }m_1,m_2\in\mathcal M .
\]
Without identification, the target is
\[
 \Theta(P;\mathcal M)=\{\theta(m):m\in\mathcal M,\ P_m=P\}.
\]
Writing \(\theta(P)\) before establishing identification can hide incompatible latent models. Confidence sets additionally need a sampling law, confidence level, asymptotic sequence and uniformity class. Point identification, coverage and informative precision are separate properties.

\textbf{Mediation.} Total effects of randomized assignment are distinct from cross-world natural indirect effects. Belief/effort-channel effects require a meaningful mediator intervention and additional measurement, exclusion or mediation assumptions. Randomization of the initial policy alone does not supply those assumptions. Report bounds or interventional alternatives when the stronger target is unsupported.

\textbf{Equilibrium.} A counterfactual \(W(z)\), \(p(z)\), or \(\mu_t^z\) requires individual optimization, feasibility, government/resource budgets, market clearing and state-transition laws. State equilibrium existence and selection, or report ranges over compatible equilibria. Initial-state integration includes subsequent endogenous transitions; current-state integration must use the policy-dependent distribution. Reduced-form invariance after policy is not assumed.

\textbf{Welfare accounts.} Scores, utility, health, dollars and ecological quantities require explicit conversion before addition. Fees, taxes, subsidies, wages and extortion payments are transfers between parties; distributional weights can make them welfare relevant, but their face value is not automatically a resource loss. State ownership, geographic boundaries, foreign-income weights and fiscal finance. Do not add profits to household welfare already including dividends, or count land capitalization and the underlying benefit twice. A benefit-cost expression must distinguish gross benefits from net welfare and subtract every real cost exactly once.

\textbf{Derivatives and risk.} Log derivatives require positive arguments; local responses require admissible perturbations and specified equilibrium branches. Differentiation under expectations needs regularity/domination, and boundaries may allow only one-sided derivatives. Risk uses a specified law; ambiguity uses a declared nonempty law set on a common history space. Adaptive policies must be nonanticipative. A robust criterion is an editorial choice unless attributed explicitly.

\textbf{Scope overlaps.} Allocation, collective choice, contracts, household bargaining and coordinated policy overlap economics with optimization and game theory. They are retained because they appeared in the original catalogue; no claim of a strictly game-theory-free selection is made.
% END ECONOMICS STATEMENT
\end{document}
